\documentclass[11pt]{article}
\usepackage[utf8]{inputenc}
\usepackage{amsmath, amssymb}
\usepackage{geometry}
\geometry{margin=1in}

\title{Summary of the Lagos (2006) Model of TFP}
\author{Gemini CLI Extraction}
\date{\today}

\begin{document}

\maketitle

\section{Introduction}
The Lagos (2006) model analyzes Total Factor Productivity (TFP) in a framework with search frictions, capital, and endogenous labor hoarding. The model features risk-neutral workers and firms, with free entry of firms.

\section{Environment and Technology}
The economy consists of a unit mass of workers and an endogenous number of firms. Each firm has a single job. Capital $k$ is rented at a flow rate $c$.

\subsection{Technology}
A match between a firm and a worker produces output using match-specific productivity $x$, hours $n$, and capital $k$:
\begin{equation}
f(x, n, k) = x \min(n, k)
\end{equation}
Capital $k$ represents capacity. Output is linear in hours up to capacity $k$.

\subsection{Productivity Shocks}
Match productivity $x$ changes according to a Poisson process with arrival rate $\lambda$. Upon a change, a new productivity $x'$ is drawn from a fixed distribution $G(\cdot)$. There is also an exogenous separation rate $\delta$.

\section{Search and Matching}
Frictions are represented by a constant-returns-to-scale matching function $m(u, v)$, where $u$ is the number of unemployed workers and $v$ is the number of vacancies.
\begin{itemize}
    \item Market tightness: $\theta = v/u$.
    \item Vacancy contact rate: $q(\theta) = m(1/\theta, 1)$.
    \item Worker contact rate: $\theta q(\theta)$.
\end{itemize}

\section{Laws of Motion and Steady State}
The time path of $(1 - u_t)H_t(x)$, the mass of matches producing at productivities $x$ or lower at time $t$, is given by:
\begin{multline}
\frac{d}{dt} [(1 - u_t)H_t(x)] = \lambda(1 - u_t)[1 - H_t(x)][G(x) - G(R_t)] + \theta q(\theta)u_t[G(x) - G(R_t)] \\
- \lambda(1 - u_t)H_t(x)G(R_t) - \lambda(1 - u_t)H_t(x)[1 - G(x)] - \delta(1 - u_t)H_t(x)
\end{multline}
The first term accounts for matches with productivities above $x$ that receive innovations below $x$ but above $R_t$. The second term represents newly formed matches. The third and fourth terms represent matches in the interval $[R_t, x]$ that get shocks below $R_t$ or above $x$. The last term accounts for exogenous separations.

\subsection{Unemployment}
In steady state, the unemployment rate $u$ is given by:
\begin{equation}
u = \frac{\delta + \lambda G(R)}{\delta + \lambda G(R) + \theta q(\theta)[1 - G(R)]}
\end{equation}

\subsection{Productivity Distribution}
The cross-sectional distribution of productivity $H(x)$ for active matches ($x \ge R$) is:
\begin{equation}
H(x) = \frac{G(x) - G(R)}{1 - G(R)}
\end{equation}

\section{Asset Values}
Let $r$ be the discount rate and $b$ the flow income of the unemployed.

\subsection{Worker Values}
The value of unemployment $U$ and the value of employment $W(x)$ satisfy:
\begin{equation}
rU = b + \theta q(\theta) \int \max[W(z) - U, 0] dG(z)
\end{equation}
\begin{equation}
rW(x) = w(x) + \lambda \int \max[W(z) - U, 0] dG(z) - (\delta + \lambda)[W(x) - U]
\end{equation}
where $w(x)$ is the wage.

\subsection{Firm Values}
The value of a vacancy $V$ and the value of a filled job $J(x)$ satisfy:
\begin{equation}
rV = -ck + q(\theta) \int \max[J(z) - V, 0] dG(z)
\end{equation}
\begin{equation}
rJ(x) = \pi(x) + \lambda \int \max[J(z) - V, 0] dG(z) - (\delta + \lambda)[J(x) - V]
\end{equation}

\section{Flow Profits and Costs}
Flow profit $\pi(x)$ is defined as residual output:
\begin{equation}
\pi(x) = x \min(n, k) - w(x) - ck - \phi n - C(x, \phi)k
\end{equation}
where:
\begin{itemize}
    \item $ck$: Rental cost of capital.
    \item $\phi n$: Variable cost (e.g., electricity or worker disutility).
    \item $C(x, \phi)k$: Fixed cost per unit of capital.
\end{itemize}
A convenient specification for the fixed cost is:
\begin{equation}
C(x, \phi) = \max(\phi - x, 0)
\end{equation}
This specification helps model labor hoarding and avoids "flat spots" in the flow profit.

\end{document}
